\chapter{Conclusion}
\label{ch:conclusion}

\section{Summary}

In this thesis, I started in Chapter~\ref{ch:survey-sampling} with defining a new methodology to optimize representativity of study participants subject to budget constraints, leveraging the properties of digital ad campaigns to prove this works in practice, at scale, and can out-perform other leading digital recruitment methods, while providing significantly more flexibility. This contributes an effective, affordable, and flexible method of online recruitment to the literature. Additionally, I make the case that there is a further contribution underneath the surface: the ability to explicitly consider representativity as something that is quantified, based on specific study outcomes, and then optimized subject to real-world constraints. This deeper implication can have implications for external validity more broadly across a variety of studies and contexts yet to be considered.

Chapter~\ref{ch:malaria} shows a real-world contribution to the development and health economics literature that uses the representative optimization of Chapter~\ref{ch:survey-sampling} to critical effect for stratified recruitment and subgroup analysis. The results provide strong empirical evidence for the real-world and scale-up effectiveness of digital ads for combatting malaria through improved health behaviors. The article also contributes a novel methodology to improve the ecological validity of ad campaigns through a practical method of natural field experiments, following the taxonomy of Harrison and List, by leveraging the remarketing tools of ad platforms in a way that, to our knowledge, was never done before.  

Finally, Chapter~\ref{ch:bebbo} presents a framed field experiment to evaluate the impact of an app through combining online recruitment and surveying with deeplinking to app downloads to match survey users and app usage data. This contributes to the literature of parenting apps, and digital parenting interventions more broadly, in the development context. The contribution is not only to provide a lack of evidence for the effectiveness of this particular app, but also to highlight important considerations in app development and evaluation design in the real world. 

\section{Implications for Policy and Practice}

The most obvious methodological lessons from this thesis are for researchers interested in running studies online. This work has provided not just methodologies, tools, and examples, but also important considerations that should guide their design. The second set of lessons are substantive: regarding the specific digital interventions studied, with direct implications for the public-health and development programs that deploy them.


\subsection{Recruiting and running studies on ad platforms}

Researchers increasingly recruit participants through digital ad platforms, and the thesis surfaces practical considerations that anyone doing so should keep in mind. The most important follows directly from how these platforms work: an ad platform, left to its own devices, optimizes for maximal engagement at minimal cost. That objective is often counter to the researcher's objective. It is therefore essential to define, explicitly and in advance, what the study actually needs---the estimand, the target population, and the covariates relevant to it---rather than accepting whatever sample the platform's optimization happens to deliver. This discipline of definition was previously often ignored. In a digital world, however, it becomes more important precisely because the default machinery is so active, high-dimensional, and opaque. The methods of Chapter~\ref{ch:survey-sampling}, together with targeting tools such as the lookalike audiences of Chapter~\ref{ch:malaria}, are examples of the kinds of levers that can be used once the researcher has defined their targets.

An additional consideration consists of the process of identity and authentication, which comes up in all digital interaction and increasingly so in an international and agentified world where fraud is rampant. The methodologies of Chapter~\ref{ch:survey-sampling}, used in both of the other chapters, keep users entirely within a platform to maintain identity, outsourcing the problem to the platform itself and minimizing fraud, in practice, at a low cost. 

\subsection{Digital ad campaigns for health behavior}

Chapter~\ref{ch:malaria} reports the first evidence, to my knowledge, of a digital ad campaign's effect on malaria-related outcomes. It shows, through the feed experiment, that the persuasive effect of the ad content holds across the full at-risk population rather than only the segment the at-scale campaign happened to reach. This matters for policy because public-health communication on social media has overwhelmingly been judged by engagement proxies---clicks, impressions, reach---rather than by the behaviors it exists to change. The lesson generalizes well beyond malaria, to other health and public behavior questions: agencies should, and now affordably can, measure behavioral outcomes directly rather than settle for engagement, and they should treat persuasive content and effective targeting as two distinct problems. A campaign may carry genuinely persuasive content and still fail at the population level if platform delivery, optimizing for cheap impressions, never reaches those most at risk---exactly the pattern in Chapter~\ref{ch:malaria}, where the ad content moved behaviour across the full at-risk population in the feed experiment, yet the at-scale campaign's detectable effects were concentrated among higher-SES, more-connected households---the lower-risk groups that engagement-optimised delivery reaches most cheaply.

\subsection{Apps and the decision to scale}

The null result of Chapter~\ref{ch:bebbo}, driven by both a lack of engagement and measurement artifacts, carries a distinct, cautionary lesson for the funders of digital interventions. The private-sector logic of apps does not transfer to public objectives: a commercial app is viable when a small, loyal minority engages deeply, whereas a public-health or development app must engage a broad share of its target population to move a population-level outcome---and most apps, Bebbo included, shed the overwhelming majority of their users within a day. Before committing budget to scale an app, funders should therefore demand evidence of broad, sustained engagement rather than intense engagement by a few, and should define precisely whom the app is meant to help. With a concrete target in mind, a study must be devised to both be representative of that target population and present the intervention in an ecologically valid context. Part of why the Bebbo evaluation found nothing is that the population it could reach was already doing well; an intervention aimed at a population with little room to improve cannot demonstrate impact however good it is. Defining the target population well is thus not only a recruitment concern but a precondition for any program that seeks to use an app as an intervention at all.


\section{Limitations of this Research}

The first chapter develops and validates a method for assembling representative survey samples by adaptively allocating advertising spend across population strata. Two evidentiary gaps bound the strength of that validation as it currently stands. The first is the absence of a head-to-head comparison against naive, un-optimized online recruitment: the chapter shows that the method matches gold-standard benchmarks, but it does not yet isolate how much of that performance is owed to the adaptive optimization itself rather than to digital ad recruitment in general. The second is that this benchmarking has so far been carried out only for the United States; the deployments in other countries establish feasibility, reach, and cost, but not accuracy against a national gold standard. Until that cross-country replication is in hand, the claim of global representativity rests on weaker evidence than the claim made for the United States. Both of these major limitations are being addressed in a forthcoming revision. 


The second chapter evaluates a Facebook campaign promoting malaria prevention in India, pairing a district-level cluster-randomized trial of the at-scale campaign with an individually-randomized feed experiment that tests the ad content itself. Several limitations qualify its findings. The eighty study districts were selected for convenience---they were the cheapest in which to recruit on Facebook---so the study population excludes the less-connected and less-populous districts and does not reflect the state as a whole. In the cluster-randomized design, individual ad exposure cannot be observed, so those estimates are intent-to-treat at the district level; the feed experiment was designed precisely to circumvent this, but only for the re-contacted sample. Budget constraints allowed recruitment to be stratified by geography and dwelling type but not balanced on other characteristics, leaving the sample heavily skewed toward young men. While we ask the young men about their household incidence, there may be a bias in their reporting which limits the external, if not the internal, validity. And the medical and heterogeneity results---the suggested post-campaign reduction in self-reported malaria incidence among higher-SES households, the parallel decline in recorded urban incidence in the administrative data, and the effects on treatment-seeking---are suggestive but underpowered, and should be read as such. The corroborative evidence from administrative data lines up extremely well with the survey data, but there is not a direct match between the campaign's geographic units and the geographic catchment areas. Similarly, the household-level socioeconomic measure used in the survey has no direct counterpart in the administrative data, which is broken down only into urban and rural facilities. While the evidence lines up very nicely, there might always be another explanation that we never thought of. 


The third chapter evaluates UNICEF's Bebbo parenting app in Serbia and Bulgaria through a randomized encouragement design, in which treated participants are offered the app and their subsequent usage is matched to survey responses. Its findings carry several limitations. The sample is one of online volunteers: its app-usage findings were validated against analytics from the general population of downloaders, but its survey-based outcomes were not. Most respondents already scored well at baseline---a pronounced ceiling effect, especially on knowledge---so there was little room to detect improvement. The online format forced a short instrument that precluded richer established scales, and it produced high attrition, though without evidence of selective attrition across arms and with adequate statistical power retained. The baseline survey itself appears to have raised knowledge in both arms, a priming effect that cannot be cleanly separated from the passage of time. Most importantly, app engagement was so low that the null result cannot distinguish an ineffective app from an unused one; and because the control group received an informational website and a substantial share of participants were already familiar with Bebbo, the results speak to the effect of encouraging new users in a context where many already know the app, rather than to the effect of the app in general. A better study design might have first observed app usage data in the wild before performing statistical power calculations, rather than performing them on theoretical usage expectations. 


\section{Future Research}

\paragraph{Adaptive representativity.} The most immediate extension of this thesis carries forward the adaptive machinery of Chapter~\ref{ch:survey-sampling}, which has already proven its value in enabling recruitment for over a dozen studies in equally many countries.  The extension will relax in turn each of the simplifying assumptions that made the original algorithm tractable, unlocking new use cases and value for researchers. 

A possible next direction treats the choice of stratification variables as an online variable-selection problem: rather than fixing the partitioning covariates in advance, one would learn in real time which covariates most efficiently partition the population to reduce estimator variance. This adapts some of the ideas from the adaptive experimentation literature \citep{Kasy2021, Caria2020, Atan2019} to the realm of survey data collection, building in particular on methods for the optimal coarse segmentation of a population \citep{Zhang2022}. The natural high-value case here is election or consumer demand forecasting, where the covariates that matter are both decisive and liable to shift over time, and where a static stratification will underperform empirically. A second direction abandons the assumption of a known, static per-stratum price in favor of learning the price function itself---price as a function of budget, estimated with explicit uncertainty---so that allocation decisions internalize the dynamic response of recruitment costs to daily spend. Alongside these, recruiting across multiple platforms, such as TikTok and YouTube rather than a single ecosystem, would mitigate platform-specific bias. Of course, it's important not to overstate this last point: multi-platform recruitment broadens the frame but cannot resolve the deeper limitation that the sampling frame remains the population of ad-reachable users who volunteer to engage.

Finally, one can easily imagine the use of LLM agents for interviewing, using new AI technology to innovate on both the survey medium (voice or chat) as well as the ability to quantify previously qualitative data (conversation, open opinion). It's easy to image that an adaptive sampling framework is precisely what an autonomous agent would need to autonomously explore the world and gather information about a particular question from humans, in the most efficient manner. 

\paragraph{Reusable frameworks for evaluating ad campaigns.} A second thread generalizes the feed experiment of Chapter~\ref{ch:malaria}---a natural field experiment that delivers real advertisements in real feeds and measures their effect on behavior through a channel unconnected to the advertiser---from a one-off design into a reusable methodology. This is no longer purely prospective. Since the malaria study, I have applied the same feed-experiment methodology to run analogous natural field experiments spanning behaviors as varied as vaccine hesitancy in Nigeria, sexual and reproductive health in Kenya, early-childhood nutrition in Indonesia, and vaping and smoking among teenagers in the United States. That a single design transfers across such heterogeneous outcomes, populations, and regulatory settings is itself evidence that it can serve as a standard, affordable instrument for evaluating public-good ad campaigns with a high and qualifiable degree of ecological validity, rather than a bespoke undertaking reinvented for each study. The natural next step is to continue to polish the methodology to be more reliable, expanding to include more platforms, and generating a reusable framework that development and policy institutions can adopt directly and easily. 

\paragraph{Quantifying and comparing study validity.} A broader, more meta-scientific direction follows from what I take to be the central conceptual move of the thesis: turning representativity from a binary property, asserted or contested in a seminar room, into a continuous quantity that can be optimized subject to a budget. The same move makes external validity comparable across studies on a common, quantified scale, addressing a need that remains largely unmet in the social sciences. Quantification permits honest, budget-constrained statements about what a given study can and cannot support, displacing binary yes-or-no debates, and it should in turn improve replicability and the credibility of decisions to scale interventions up.

One can imagine, speculatively, a further step along this path. The thesis optimizes a metric at the point of recruitment; one might instead seek a statistic that validates an arbitrary sample for an arbitrary conclusion and target population, computed after the fact rather than designed for in advance. Any such statistic would be inherently relative to the estimand and the target: a sample that is representative for one question need not be representative for another. I raise this only as a possibility worth imagining, not one in which I have done any preparatory work.  

\paragraph{Institutional scale and demand.} I close on a more concrete and encouraging note. The methodologies developed here are not awaiting a hypothetical audience: they have powered and are powering over a dozen active randomized controlled trials and surveys in 2025 and 2026, with groups including the World Bank. This demonstrated demand points to a real opportunity to run studies at scale, under explicit budget constraints and with a high degree of ecological validity, in direct service of the needs of policymakers and researchers. That, ultimately, is the unifying contribution I hope this thesis makes: a small addition in what will necessarily be a large literature of methodologies that allow researchers to consider external validity, not merely as conceptual ideals but instead as practical, quantified, and honest features of empirical work subject and optimized according to real-world constraints, and available at scale. 


\printbibliography[heading=subbibliography, title={References for Chapter \thechapter}]
